\scanpage{119}
\mspage{106}
On a alors
\[
 X_U=u_!(X)(1),\qquad X_I=u_!(X)(0)
\]
et le morphisme de transition
$X_U\longrightarrow X_I$ est déduit de \inserted{la flèche} $0\xrightarrow[\unsure{\delta|d}]{}1$ dans
$\Delta^1$, et de $=$ pour $Y$. Soit
\[
 \unsure{\xi|\zeta}=u_!(X),\qquad \eta=u_!(Y)
\]
donc le diagramme (6.2.1) s'identifie
au diagramme déduit de
\[
 \varphi\colon \xi\longrightarrow\eta,\qquad \varphi=u_!(f),
\]
comme

\diagram{\begin{tikzcd}[ampersand replacement=\&,column sep=large]
\& \xi(1) \arrow[r,"{\xi(\delta)}"] \arrow[d,"{\varphi(1)}"'] \arrow[dr,phantom,"{\unsure{\text{cocart}|\text{cart}}}"] \& \xi(0) \arrow[d,"{\deleted{\unsure{\varphi(0)}}}"] \arrow[dd,bend left=50,preaction={draw,dash pattern=on .5pt off 6pt,line width=1.1pt,-},"{\varphi(0)}"]\\
(6.2.3) \& \eta(1) \arrow[r,"{\deleted{\unsure{\eta(\delta)}}}"] \arrow[dr,dashed,"{\eta(\delta)}"'] \& \widetilde{\eta}_{\unsure{1}}\deleted{\unsure{\eta_{1,0}}} \arrow[d,"{\widetilde{\varphi}\deleted{\unsure{\varphi_{1,0}}}}"']\\
\& \& \unsure{i}\eta(0)
\end{tikzcd}}

Or par 6.1.10 appliqué à $u\colon I\unsure{\text{---}|\longrightarrow}\Delta^1$,
la \unsure{relation|condition} $f\in C(I)$ resp $f\in TC(I)$
implique $\varphi\in C(\Delta^1)$ resp. $\varphi\in TC(\Delta^1)$.
Donc pour établir 6.1, on est ramené
au cas particulier où $I=\Delta^1$, $U=\{1\}$, $V=I$,
donc à prouver le

\uline{Lemme 6.2.4.} Soit $\varphi\colon\xi\longrightarrow\eta$ une
flèche dans $M(\unsure{\Delta^1|\Delta})$. Si $\varphi\in C(\Delta^1)$ (resp.
% archive sheet number: 118

